% Generado por Experimentos/generar_resultados_latex.py desde complementario.
% No editar a mano: volver a ejecutar el script.

\begin{table}[ht]
\centering
\footnotesize
\begin{tabular}{c c c c}
\hline
Ciclo & $\alpha=0$ & $\alpha=0.5$ & $\alpha=1$ \\
\hline
0 ($M_0$) & \multicolumn{3}{c}{$0.9896$ [$0.9841$; $0.9933$]} \\
1 & $0.9839$ [$0.9813$; $0.9877$] & $0.9839$ [$0.9817$; $0.9849$] & $0.9831$ [$0.9802$; $0.9853$] \\
2 & $0.9837$ [$0.9786$; $0.9865$] & $0.9865$ [$0.9802$; $0.9909$] & $0.9721$ [$0.9663$; $0.9782$] \\
3 & $0.9893$ [$0.9853$; $0.9925$] & $0.9856$ [$0.9806$; $0.9909$] & $0.9532$ [$0.9508$; $0.9567$] \\
4 & $0.9877$ [$0.9833$; $0.9913$] & $0.9849$ [$0.9810$; $0.9881$] & $0.9529$ [$0.9484$; $0.9575$] \\
5 & $0.9854$ [$0.9837$; $0.9877$] & $0.9839$ [$0.9778$; $0.9877$] & $0.9401$ [$0.9357$; $0.9476$] \\
\hline
\end{tabular}
\caption{Distinct-3: media y rango [mínimo; máximo] de las tres cadenas por ciclo y proporción sintética $\alpha$. $M_0$ es compartido por las tres proporciones dentro de cada repetición. Salvo la perplejidad, calculada con los 128 contextos de prueba, las métricas corresponden al análisis complementario \textit{post hoc} con los 90 contextos en los que todas las etapas alcanzan $W=30$ palabras.}
\label{tab:medias_distinct_3}
\end{table}
